\section{\ref{sec:muscles}장. 근육으로의 출력}

출력의 구조, 곧 운동 단위, 시냅스의 안전 여유, 크기 순서대로의 동원, 근육 쌍, 신장 루프, 리듬 발생기는 직접 측정되었다. 지연 루프의 안정 조건은 정리이다. 몸의 내부 모델은 근거가 탄탄한 가설이다. 그림의 수치는 이 책의 모델에 따른 계산이다.

{\footnotesize
\setlength{\tabcolsep}{3pt}\renewcommand{\arraystretch}{1.15}
\begin{longtable}{>{\raggedright}p{0.5cm}>{\raggedright}p{4.0cm}>{\raggedright}p{5.6cm}>{\raggedright}p{2.3cm}>{\raggedright\arraybackslash}p{1.7cm}}
\toprule
No. & 명제 & 실험과 측정한 것 & 출처 & 상태 \\
\midrule\endfirsthead
\toprule
No. & 명제 & 실험과 측정한 것 & 출처 & 상태 \\
\midrule\endhead
1 & 운동 단위: \nt{ACET} 뉴런 하나가 수백에서 이천 개 남짓의 섬유 무리를 제어한다; 섬세한 조정을 하는 근육에서는 단위가 더 작다 & 사람 근육, 운동 축삭과 근육 섬유 계수: 뉴런당 섬유는 평균적으로 손의 골간근에서 약 $340$개, 완요골근에서 약 $410$개, 안쪽 장딴지근에서 약 $1930$개. 가장 작은 단위의 정확한 수는 본문에 없다. & \cite{Feinstein1955mu} & 측정됨 \\
2 & 근육 위의 시냅스는 섬유를 동작시키는 데 필요한 양의 몇 배나 되는 신경전달물질을 방출한다 & 신경근 시냅스 측정 총설: 성체 포유류의 안전 여유(방출량 대 문턱 양의 비)는 보통 $3$--$5$; 사람에서는 여유가 방출량보다 수신 측 막의 주름에 더 많이 기댄다. & \cite{WoodSlater2001} & 측정됨 \\
3 & 연축~\eqref{eq:twitch}은 $50\ms$ 정도의 시간 $T_c$ 동안 커진다 & 수의적 힘을 낼 때 사람 손의 단일 운동 단위, 단위의 스파이크로 평균한 힘: 상승 시간 $30$--$100\ms$, 단위의 $80\,\%$가 $70\ms$ 미만. 함수 $(t/T_c)e^{1-t/T_c}$는 단위 풀 모델의 근사. & \cite{MilnerBrown1973rec, Fuglevand1993} & 측정됨 \\
4 & 연축의 합~\eqref{eq:forcesum}: $5$, $15$, $40$ Hz에서 힘 $0.2$--$1.1F_1$, $1.8$--$2.2F_1$, 약 $5.4F_1$(그림~\ref{fig:tetanus}) & \texttt{models/\allowbreak muscle\_\allowbreak models.py}에 따른 계산, $T_c=50\ms$: 마지막 $0.3\,\text{s}$에서 $0.21$--$1.10$, $1.76$--$2.19$, $5.28$--$5.49\,F_1$. 선형 합은 단순화: 실제 근육의 포화는 모델에 없다. & \texttt{models/\allowbreak muscle\_\allowbreak models.py} & 모델 \\
5 & 단위는 크기 순서대로 켜진다; 가장 약한 단위는 가장 강한 단위보다 백 배 약하다 & 고양이 배쪽 뿌리: 반사 입력이 커질 때 뿌리에서 스파이크가 작은 뉴런, 곧 작은 뉴런이 먼저 방전한다. 사람 손의 골간근: 단위의 연축 힘은 $0.1$에서 $10$~g이고, 단위가 켜지는 힘에 거의 선형으로 커진다. & \cite{Henneman1957, MilnerBrown1973rec} & 측정됨 \\
6 & 힘의 단계는 힘에 비례한다, $\Delta F/F\approx q-1$~\eqref{eq:sizeprinciple}: 부동소수점 & 힘의 등비수열에서 본문이 유도. 실험과 부합: 큰 힘에서는 같은 힘 증가에 새로 켜지는 단위가 훨씬 적다. & 이 장의 유도; \cite{MilnerBrown1973rec} & 이론 \\
7 & 힘의 퍼짐은 힘 자체에 비례해 커진다 & 엄지 근육의 수의적 등척성 힘: 힘의 표준편차는 평균 힘에 따라 선형으로 커진다; 전기 자극으로 같은 힘을 낼 때는 퍼짐이 일정하다. 단위 풀 모델: 힘 순서대로의 동원이 선형 증가를 준다. & \cite{Jones2002sdn} & 측정됨 \\
8 & 움직임의 매끄러움은 신호 의존 잡음의 결과이다 & 모델: 이런 잡음에서 끝점 퍼짐을 최소화하는 궤적이 측정된 눈과 팔의 궤적을 재현한다. 이 원인을 다른 원인과 구별하는 직접 실험은 없다. & \cite{HarrisWolpert1998} & 가설 \\
9 & 근육 쌍의 힘 차가 회전력을, 합이 강성을 정한다~\eqref{eq:antagonists} & 공동 긴장에 의한 임피던스 제어 이론. 사람 팔: 작은 밀기로 잰 각 관절의 강성은 그 관절의 회전력에 대략 비례한다; 공동 긴장의 비율이 다르면 손 강성 타원의 모양과 방향이 바뀐다. & \cite{Hogan1984, GomiOsu1998stiff} & 측정됨 \\
10 & 신장 센서는 자기 \nt{ACET} 뉴런을 흥분시키고, 억제성 중개 뉴런을 거쳐 길항근을 끈다(그림~\ref{fig:antagonists}) & 고양이 척수: 신장 센서 섬유로 흥분되는 단일 중개 뉴런이 길항근 \nt{ACET} 뉴런에 단일 시냅스 억제 전위를 만든다, 시냅스 지연 $0.28$--$0.42\ms$. 중개 뉴런의 신경전달물질(글리신)은 이 실험에서 확인하지 않았다. & \cite{Jankowska1972Ia} & 측정됨 \\
11 & 루프 지연: 척수 루프 $25$--$30\ms$, 피질을 거치면 1.5--3배, 눈을 거치면 $100$--$200\ms$ & 사람 팔, 관절 밀기: 근육의 짧은 응답은 $20$--$45\ms$, 긴 응답은 $45$--$105\ms$, 수의적 응답은 $120$--$180\ms$ 뒤. 움직이는 동안 보이는 손가락 위치를 옮기면: 보정은 평균 $160\ms$ 뒤. & \cite{Pruszynski2008rapid, SaundersKnill2003vis} & 측정됨 \\
12 & $\dd x/\dd t=-Gx(t-d)$는 $Gd<\pi/2$일 때 안정하다~\eqref{eq:delaystable}; 경계에서 주파수는 $1/(4d)$ & 지연 방정식의 근에 관한 정리. \texttt{models/\allowbreak muscle\_\allowbreak models.py}로 확인, $d=30\ms$: $3\,\text{s}$까지의 진폭은 초당 $G=40$에서 $3\cdot10^{-6}$, $G=50$에서 $0.13$, $G=55$에서 $38$(경계 $52$). & \cite{Hayes1950} & 이론 \\
13 & 생리적 떨림 $8$--$12$~Hz; 신장 루프는 그 원천 가운데 하나 & 측정 총설: 약 $10$~Hz 대역의 미세 떨림은 건강한 사람에게 있다; 그 기원은 복합적이다. 중추 진동, 단위 방전의 성질, 기계적 공진, 반사 루프의 공진이 있고, 조건에 따라 우세한 것이 다르다. & \cite{McAuleyMarsden2000} & 가설 \\
14 & 뇌는 몸의 내부 모델로 명령의 결과를 예측한다 & 사람이 물체를 손가락으로 집어 들고 팔을 움직인다: 관성, 점성, 탄성 부하에서 쥐는 힘이 루프 지연 없이 부하와 동시에 변한다. 곧 부하가 예측된다. 총설이 이런 실험을 정리한다; 뉴런에서 모델 자체를 직접 관찰한 것은 없다. & \cite{FlanaganWing1997grip, WolpertGhahramani2000} & 가설 \\
15 & 걷기, 숨쉬기, 씹기의 리듬은 일정한 입력에서 국소 네트워크가 만든다 & 실험 총설: 분리된 신경계(척수, 신경절)는 리듬적 입력도 근육의 되먹임도 없이 리듬 운동 패턴을 낸다. & \cite{MarderBucher2001} & 측정됨 \\
16 & 피로를 동반한 상호 억제~\eqref{eq:halfcentre}는 주기 $0.84\,\text{s}\approx2\tau_a$의 리듬을 준다(그림~\ref{fig:halfcentre}) & 정리: 상호 억제와 적응을 가지며 일정한 입력에서 안정 상태가 없는 네트워크는 반드시 진동한다. \texttt{models/\allowbreak muscle\_\allowbreak models.py}에 따른 계산($I=1$, $\beta=2$, $b=2$, $\tau=20\ms$, $\tau_a=0.4\,\text{s}$): 주기 $0.84\,\text{s}$. 실제 발생기가 바로 이렇게 만들어져 있다는 것은 보이지 않았다. & \cite{Matsuoka1985osc} & 모델 \\
\bottomrule
\end{longtable}}
